% !PRINTED 233
\weberpage{233}
%<<CONT
beschränken uns bei diesem Beweise auf die Betrachtung der Ikosa\"edergruppe, da für die anderen Gruppen ganz ähnliche Schlüsse zu machen sind, die wir dem Leser überlassen können. Der Umstand, dass bei der Tetra\"eder- und Octa\"edergruppe neben den absoluten auch relative Invarianten vorkommen, während die Ikosa\"edergruppe als einfache Gruppe nur absolute Invarianten hat, ist hierbei von keinem wesentlichen Einflusse.

Wir beweisen also, dass alle Invarianten der Ikosa\"edergruppe sich als ganze rationale Functionen der drei Formen $f$, $H$, $T$ darstellen lassen.

Dazu führt uns folgende Schlusskette:

1. \gesperrt{Keine zwei der Ikosa\"ederformen $f$, $H$, $T$ haben einen gemeinschaftlichen Theiler.}

Dies folgt unmittelbar aus der Definition dieser Functionen, wonach die Gleichungen $f = 0$, $H = 0$, $T = 0$ die fünfzähligen, dreizähligen und zweizähligen Pole liefern, die alle von einander verschieden sind.

2. \gesperrt{Eine Doppelwurzel der Ikosa\"edergleichung}
\weq{1}{H^3 - z f^5 = 0}
\gesperrt{ist nothwendig eine Wurzel von $f$, $H$ oder $T$, und kann also nur für einen der Werthe $z = 0, \infty, 1728$ eintreten.}

Wenn nämlich (1) eine Doppelwurzel hat, so müssen mit der Function zugleich die erste Ableitung, oder, wenn man die homogene Form anwendet, nach der \gesperrt{Euler}'schen Theorie [Bd. I, \S.\,17, (5)] die beiden Ableitungen nach $x_1$ und $x_2$ zugleich verschwinden, also
\[
3H^2 H'(x_1) - 5z f^4 f'(x_1) = 0,
\]
\[
3H^2 H'(x_2) - 5z f^4 f'(x_2) = 0;
\]
folglich bleiben, da $H$ und $f$ nach 1. nicht zugleich verschwinden, nur drei Möglichkeiten: entweder $H = 0$, $z = 0$, oder $f = 0$, $z = \infty$, oder endlich
\[
H'(x_1)\,f'(x_2) - H'(x_2)\,f'(x_1) = 0,
\]
d. h. $T = 0$, $z = 1728$ [\S.\,60, (4)]. Hieran schliesst sich auch der evidente Satz:
%CONT>>
